\section{Representation and Editability}
\label{sec:representation}
We treat the \emph{delivered artifact}, rather than the architecture of its generator, as the first classification axis.
In design production the image rendered to viewers and the editable artifact opened by designers are distinct observables.
Table~\ref{tab:representation} summarizes common affordances and the evidence that is still needed.

% Independent source table.
\begin{table*}[t]
\centering
\caption{Representation--editability crosswalk.
The listed operations are possible affordances, not guarantees of empirically validated editing.}
\label{tab:representation}
\small
\begin{tabularx}{\textwidth}{@{}p{0.18\textwidth}XX@{}}
\toprule
\textbf{Delivered artifact} & \textbf{Potential post-generation operation} & \textbf{Critical evidence boundary} \\
\midrule
Layout boxes and classes & Move, resize, or constrain element geometry. & Geometry does not establish editable text, asset identities, or a portable design file. \\
Flattened raster image & Pixel-level modification, inpainting, or mask-based regeneration. & Raster image editing is valid, but it does not imply native vector or text-object editing. \\
SVG vector paths & Modify paths, fill, stroke, and layer order; interpolate shapes. & Outlined letters are not semantic font objects; path access is not an editing-success score. \\
Layered image/document & Reposition, replace, or re-render separable layers and components. & RGBA layers may still contain fused objects; alpha channels alone do not ensure semantic separation. \\
Code-native HTML/CSS or presentation objects & Edit DOM/object nodes, text, and layout rules; re-render and inspect. & Syntactically valid code need not preserve intended layout or succeed at natural-language edit requests. \\
\bottomrule
\end{tabularx}
\end{table*}


\subsection{Layout Geometry and Raster Composites}
\ac{layoutdm} models category and quantized spatial attributes through discrete diffusion with conditional generation from partial layouts.
\ac{layoutflow} models continuous geometric generation.
Such systems provide useful structural arrangements, but layout coordinates need not include the image and text assets needed for a fully editable poster.
Content-aware layout systems additionally place foreground elements in visual context \citep{cglgan2022,layoutdetr2024}.

Raster-oriented design generation produces integrated image renderings, including diffusion-based design images \citep{designdiffusion2025}, product posters with rendered typography \citep{postermaker2025}, and images with visually generated text \citep{textdiffuser22024}.
A \ac{png} image can retain an alpha channel, but the file format alone does not preserve independently editable text or design objects.
Visual quality and text-rendering accuracy are valuable outcomes, but they do not measure whether a later requested edit preserves unrelated regions or text objects.
SAWNA exemplifies a boundary: it enforces user-reserved blank regions during image synthesis to support later composition \citep{sawna2025}.
Those regions are useful but do not establish post-export semantic object manipulation.
\ac{postertext} targets localized rendered text patches; its \ac{html} structure is used to construct annotated training pairs, not exported as a native editable document by the generation model (Section 3 and Figure 3).
Quantitative raster editing need not preserve independently manipulable original font, text or vector objects.

\paragraph{Document structures and conditioned layouts.}
READ uses a recursive latent tree to capture relationships among document elements \citep{read2020}, whereas Neural Design Network incorporates user constraints directly in structured layout generation \citep{ndn2020}.
\ac{aclayoutgan} models the association between content attributes and layout decisions, and multi-choice layout learning encourages diversity among plausible arrangements \citep{layoutmcl2023}.
In image-aware composition, \ac{pdagan} adapts between different background domains with a pixel-level discriminator, whereas \ac{layoutdit} uses a diffusion transformer to balance background content and graphic elements \citep{layoutdit2024}.
Visual Layout Composer explicitly models a coupled image/vector diffusion process \citep{vlc2024}, illustrating how a predicted layout, composed image, and vector representation can be learned together while remaining distinct deliverables.

\subsection{Vector Path and Layer Structures}
\ac{deepsvg} models drawing commands and studies reconstruction and latent interpolation; \ac{svgformer} learns representations of continuous vector graphics.
\ac{vectorfusion} optimizes differentiable \ac{svg} paths with text-to-image diffusion priors.
\ac{neuralsvg} learns an implicit representation for layered vector synthesis, whereas \ac{omnisvg} tokenizes vector commands and coordinates for multimodal generation.
\ac{iconshop} models \ac{svg} icon paths and demonstrates editing by path-level infilling.
Layer-wise image vectorization instead aims to infer multiple stylized regions and corresponding vector shapes from an input raster \citep{live2022}.
The existence of vector paths does not establish instruction-driven edit accuracy or unchanged-object preservation.

Layered design synthesis and raster-to-layer reconstruction should be distinguished because they can expose similar compositing operations while solving different inference problems.
Hierarchical systems such as \ac{cole} and \ac{opencole} construct layered designs through intermediate composition decisions.

Raster-to-layer recovery begins with a finished rendering rather than a design brief.
\ac{layerd} reconstructs layers from flattened designs, while \ac{qwenimagelayered} predicts a variable number of raster \ac{rgba} layers \citep{layerd2025,qwenimagelayered2026}.
\ac{pixelart} explores pixel-space rectified-flow decomposition with additional layout conditioning \citep{jia2026pixelart}.
Although these outputs can expose separable content for recomposition, reconstruction metrics and qualitative edits do not establish native text-object preservation or reliable semantic revision.

Transparent-asset generation provides compositing control even when a design-object hierarchy is absent.
\ac{layerfusion} instead jointly synthesizes transparent foreground and background layers, while \ac{igd} combines predicted layer attributes with generated image assets.
\ac{alfie} instead obtains a transparent raster foreground by estimating an alpha matte from diffusion-transformer attention maps, enabling placement in an external composition but not independent editing of internal text or vector objects.
A portable \ac{rgba} illustration is therefore a compositing affordance, not evidence of a native multi-object design hierarchy.

Recovering reusable layers from a raster composite is not equivalent to recovering the original document hierarchy.
\ac{layeringdiff} separates a composite into foreground RGB and alpha, reconstructed background RGB, and a recomposed image (Section 3 and Figure 1).
Its iterative examples (Figure 7) demonstrate further foreground separation but not recovery of original font or vector objects.
Consequently, a reconstruction evaluation should not be interpreted as proof of native semantic editing.

Instruction-conditioned recovery can also target ordered, partially occluded design components.
\ac{uniworlddesign} generates standalone transparent assets and recovers ordered, complete semantic \ac{rgba} layers from a flattened image using global and per-layer instructions (Sections 3.3--3.5).
Although designer-authored \ac{psd} files supply its layer-tree supervision, the reported model output is an \ac{rgba} stack, not an exported editable PSD document with native font, vector, or typography objects.
A complete occluded raster object may be moved and recomposed, but that format affordance is not itself a measured success rate for post-export semantic design edits.

Code-native designs afford a different form of control because text and layout can be addressed through document structure rather than image regions.
Systems that deliver \ac{html}/\ac{css} expose individually manipulable elements and properties \citep{designascode2026,editablevisual2026}.
This source-level access broadens the possible editing operations, but only a controlled edit protocol can establish whether requested changes succeed without collateral damage.

\subsection{Four Types of Evidence}
We distinguish: \emph{format affordance}, where structure permits a type of manipulation; \emph{workflow demonstration}, where authors show specific edits; \emph{human editing efficiency}, where users are measured completing a design task; and \emph{quantitative edit-operation evaluation}, where target change and unintended changes are scored.
The levels are not a universal ordinal quality scale.
A qualitative demonstration can be convincing for one operation while failing to establish statistical reliability, and a controlled human study can measure a different property from automatic instruction following.
